(1) 两平行板间距 \(b\)，流体1处于 \(x<0\)，流体2处于 \(x>0\)，黏度为 \(\mu_1,\mu_2>0\)。两者以速度 \(U>0\) 沿 \(+x\) 共同平移；板隙薄、不可压、低Re，面内速度服从 \(\boldsymbol u_j=-\kappa_j\nabla p_j\)，\(\kappa_j=b^2/(12\mu_j)\)。在以 \(U\) 平移的坐标系中令界面 \(x=\eta(y,t)\)，说明压强扰动应满足的方程，写出在无穷远衰减的正弦扰动形式。
(2) 取 \(\eta=\eta_0e^{iky+st}\)，在界面使用运动学条件及线性Laplace压差 \(p_1-p_2=-\gamma\eta_{yy}\)（\(\gamma>0\) 为界面张力），严格推导增长率 \(s(k)\)，标明界面基态压强梯度也必须计入。
(3) 指出何种黏度排序导致指进不稳定。求不稳定波数范围、最快增长波数及其增长率，检查 \(U=0\)、\(\gamma=0\) 与 \(\mu_1=\mu_2\) 的极限。
(4) 验证短波的毛细稳定项符号与量纲；给出小斜率界面表面能二阶增量，解释不稳定时能量从何而来，并讨论薄隙近似、有限板宽及扩散性界面所限制的波数范围。
